\section{Introduction}
\label{sec:intro}

Consider a pharmacy support queue. A message arrives: \emph{``I need a
refill on my thyroid tablets, and are you open on Sunday? Also I've been
getting dizzy since starting them.''} A system routing this must decide,
simultaneously, which topics it concerns (two), whether it mentions
clinical risk that a human should see (it does), and how urgent it is.
Each of these is a decision over a \emph{known, closed} set of answers.

The dominant approach is to ask a language model and parse what comes back.
This works, and it is expensive in three distinct ways. Each question costs
its own generation. The output is unconstrained, so it must be coerced onto
the menu by grammar-constrained decoding or by retrying. And the resulting
token probabilities are a poor basis for the thresholding that operational
systems need, because they describe a sampling distribution over strings
rather than a calibrated belief over options.

\subsection{Typed decisions}

We use \emph{typed decision} for a question whose answer set is fixed in
advance and known to the system. Three types cover most of what production
routing and guardrail systems need:

\begin{itemize}[leftmargin=1.4em,itemsep=0.25em]
  \item \choice{} --- exactly one of $K$ options. Intent, category, routing
    target. $K$ reaches 151 in our evaluation.
  \item \score{} --- one level on an \emph{ordered} scale. Severity,
    urgency, quality. The ordering is information a \choice{} discards.
  \item \noul{} --- yes or no. Flags and safety gates. Several asked
    together express multi-label structure, which a single \choice{}
    cannot: a softmax over intents has no way to say ``both''.
\end{itemize}

\subsection{What this paper is about}

\openjev{} is an open implementation of an architecture for answering many
such questions about one input at once, and an empirical account of what
makes it work. We set out to reproduce the behaviour of a commercial
typed-decision API, and the interesting result is not the architecture,
which is a recombination of published ideas (\S\ref{sec:background}). It is
which parts of the recipe actually carry the performance.

Three findings organise the paper.

\paragraph{The architecture does not generalize on its own.} A pretrained
bidirectional encoder, given the option-scoring output format and no task
training, scores $0.7\times$ chance on held-out schemas --- worse than
guessing. Its larger sibling reaches $2.2\times$. The output format makes
zero-shot scoring \emph{possible}; it does not make it \emph{present}
(\S\ref{sec:res-architecture}).

\paragraph{Generalization is a property of the training mixture, and some
of it is synthesizable.} Our first mixture left held-out transfer at
$0.9\times$ chance. The single largest fix was not more data but a
different \emph{shape} of data: padding training menus with distractor
labels borrowed from other tasks, which moved a 77-way held-out menu from a
literal $0.000$ to $0.205$. A model that has never seen a large menu
collapses onto one label when handed one, and you can manufacture large
menus without finding datasets that have them (\S\ref{sec:res-data}).

\paragraph{With labels, small models win; without them, they do not.} A
rank-16 adapter on a 1.7B backbone, trained on 395 examples, beats the
commercial API on three of seven router measures and ties the rest. The
same architecture without task labels remains far behind it. That split is
the honest summary of where this technology stands
(\S\ref{sec:res-router}, \S\ref{sec:discussion}).

\subsection{Contributions}

\begin{enumerate}[leftmargin=1.5em,itemsep=0.25em]
  \item An architecture and open implementation for multi-question typed
    decisions in a single forward pass, with latency flat in the number of
    questions (10 questions cost 6\% more than 1) and full-precision
    calibrated probabilities.
  \item \emph{Label-space augmentation}, a training-time intervention that
    manufactures the menu-size diversity a mixture lacks, and which is the
    difference between no transfer and useful transfer at high cardinality.
  \item A measured account of parameter-efficient fine-tuning for this task
    shape, including the result that a 17M-parameter adapter outperforms
    opening all 1{,}725M parameters, at $1/39$ the artifact size.
  \item A reproducible benchmark suite with an enforced contamination
    guard, a held-in control that distinguishes ``no transfer'' from ``a
    bug'', and a hardened routing task built to retain headroom where the
    reference system saturated.
  \item A set of negative and retracted results reported at the same
    prominence as the positive ones (\S\ref{sec:ablations}), including two
    interventions that made things worse and one conclusion we withdrew
    after finding the bug that produced it.
\end{enumerate}

Everything is reproducible from a public repository, with released
checkpoints, adapters, and datasets listed in Appendix~\ref{app:artifacts}.
